\documentclass[10pt]{article}
\usepackage[T1]{fontenc}
\usepackage{lmodern}
\usepackage[margin=22mm]{geometry}
\usepackage{xcolor}
\usepackage{microtype}
\usepackage{enumitem}
\usepackage{hyperref}
\usepackage{xurl}
\usepackage{forest}
\usetikzlibrary{arrows.meta}
\usepackage{pdflscape}
\usepackage[most]{tcolorbox}
\usepackage{titlesec}
\usepackage{fancyhdr}
\definecolor{Navy}{HTML}{17324D}
\definecolor{Teal}{HTML}{176B6B}
\definecolor{PaleBlue}{HTML}{EAF2F8}
\definecolor{FileFill}{HTML}{F7F9FA}
\definecolor{DescText}{HTML}{334E5C}
\definecolor{TreeLine}{HTML}{607D8B}
\definecolor{Warn}{HTML}{8A4B08}
\hypersetup{colorlinks=true,linkcolor=Teal,urlcolor=Teal,pdfborder={0 0 0}}
\titleformat{\section}{\Large\bfseries\color{Navy}}{}{0pt}{}[\color{Teal}\titlerule]
\titleformat{\subsection}{\large\bfseries\color{Teal}}{}{0pt}{}
\setlist[itemize]{leftmargin=*,itemsep=.45em,topsep=.45em}
\setlist[enumerate]{leftmargin=*,itemsep=.45em,topsep=.45em}
\setlength{\parindent}{0pt}
\setlength{\parskip}{.65em}
\setlength{\emergencystretch}{3em}
\newcommand{\breakableunderscore}{\textunderscore\allowbreak}
\newcommand{\code}[1]{\begingroup\ttfamily\let\_\breakableunderscore#1\endgroup}
\forestset{
 rootnode/.style={draw=Navy,fill=Navy,text=white,rounded corners=2pt,inner sep=5pt},
 filenode/.style={draw=TreeLine,fill=FileFill,rounded corners=2pt,inner sep=3.2pt},
}
\pagestyle{fancy}\fancyhf{}\fancyhead[L]{\small\color{DescText}Repository Structure Inventory}\fancyhead[R]{\small\color{DescText}\leftmark}\fancyfoot[C]{\thepage}
\setlength{\headheight}{14pt}
\newtcolorbox{verificationbox}{breakable,colback=orange!5,colframe=Warn,title={Verification status},fonttitle=\bfseries}
\begin{document}
\hypersetup{pageanchor=false}
\begin{titlepage}
\pagecolor{PaleBlue}
\color{Navy}
\vspace*{0.17\textheight}
{\Huge\bfseries Grade 11 — Term 3 Structure\par}
\vspace{1em}
{\Large Repository resource inventory\par}
\vfill
{\large Authoritative conversion of \code{STRUCTURE.md}\par}
\end{titlepage}
\nopagecolor
\hypersetup{pageanchor=true}
\tableofcontents
\clearpage

\section{Overview}

Grade 11 Term 3 is a LaTeX-based collection for teaching and assessing
hypothesis testing. It progresses from designing hypotheses (criteria C1–C3),
through one-sample proportion and mean tests and interpretation (C4–C7), to a
web implementation and report using two-sample tests (C8–C10). The four
top-level resource directories separate guided practice, learning evidence,
catch-up work, and a longer worked exam. There are no HTML, CSS, or JavaScript
implementations in this directory; the learning-evidence brief asks students to
create and submit those files.

Audience labels below distinguish \textbf{student-facing}, \textbf{teacher/solution}, and
\textbf{shared lesson} resources. LaTeX sources are printable lesson, assessment, or
answer-key documents; CSV files are data resources rather than generated build
outputs.

\section{Directory tree}

\begin{tcolorbox}[colback=PaleBlue,colframe=Teal,title={Reading the diagrams}]
The directory tree is divided into one landscape diagram per top-level directory so that every documented entry and description remains legible. Directory nodes use a dark fill; file nodes use a light fill, with descriptions shown in a secondary text color. The inventory root is \code{.}.
\end{tcolorbox}
\begin{landscape}
\subsection{\code{1-practice\_activities/}}
\begin{center}
\begin{forest}
for tree={grow=east, parent anchor=east, child anchor=west, anchor=west, edge={-Latex,draw=TreeLine}, l sep=9mm, s sep=1.2mm, font=\sffamily\fontsize{6.2}{7.2}\selectfont},
[{\textbf{1-practice\_activities/}\\\textcolor{DescText}{Shared lesson practice and reference material, mostly with worked solutions.}}, rootnode, text width=4.5cm
[{\code{G11\_T3\_L1\_C1C2C3\_practice\_activity.tex}\\\textcolor{DescText}{Shared lesson: extended C1–C3 hypothesis-design practice with ten mean/proportion contexts and solutions.}}, filenode, text width=16.5cm]
[{\code{G11\_T3\_L1\_C1C2C3\_practice\_activity1.tex}\\\textcolor{DescText}{Shared lesson: introductory C1–C3 activity with four solved contexts for parameters, hypotheses, and tail direction.}}, filenode, text width=16.5cm]
[{\code{G11\_T3\_L1\_C1C2C3\_practice\_activity2.tex}\\\textcolor{DescText}{Shared lesson: eight solved C1–C3 tests emphasizing directional claims, changing alpha, and one- versus two-tailed comparisons.}}, filenode, text width=16.5cm]
[{\code{G11\_T3\_L1\_C1C2C3\_practice\_activity3.tex}\\\textcolor{DescText}{Shared lesson: six open-ended, locally framed C1–C3 problem-design prompts with solution guidance.}}, filenode, text width=16.5cm]
[{\code{G11\_T3\_L2\_C4\_practice\_activity.tex}\\\textcolor{DescText}{Shared lesson: solved C4 ladder for one-sample proportion z-tests, from basic cases to threshold/sample-size comparisons.}}, filenode, text width=16.5cm]
[{\code{G11\_T3\_L2\_C5\_practice\_activity.tex}\\\textcolor{DescText}{Shared lesson: parallel solved C5 ladder for one-sample mean z-tests in count contexts.}}, filenode, text width=16.5cm]
[{\code{G11\_T3\_L3\_C6\_practice\_activity.tex}\\\textcolor{DescText}{Shared lesson: solved C6 cases on decisions, contextual conclusions, Type I/II errors, and policy implications.}}, filenode, text width=16.5cm]
[{\code{G11\_T3\_L3\_C7\_practice\_activity.tex}\\\textcolor{DescText}{Shared lesson: solved C7 minimum-sample-size problems with one-/two-tail and multi-alpha graphs.}}, filenode, text width=16.5cm]
[{\code{G11\_T3\_MATERIAL.tex}\\\textcolor{DescText}{Shared lesson bank: twenty scaffolded hypothesis-test contexts spanning school, consumer, and economics topics.}}, filenode, text width=16.5cm]
[{\code{G11\_T3\_summary.tex}\\\textcolor{DescText}{Teacher/reference overview summarizing the intended progression and problem roles of the C1–C7 practice sequence.}}, filenode, text width=16.5cm]
]
\end{forest}
\end{center}
\end{landscape}
\begin{landscape}
\subsection{\code{2-learning\_evidences/}}
\begin{center}
\begin{forest}
for tree={grow=east, parent anchor=east, child anchor=west, anchor=west, edge={-Latex,draw=TreeLine}, l sep=9mm, s sep=1.2mm, font=\sffamily\fontsize{6.2}{7.2}\selectfont},
[{\textbf{2-learning\_evidences/}\\\textcolor{DescText}{Student assessments, answer keys, a web-project brief, and its candidate/example datasets.}}, rootnode, text width=4.5cm
[{\code{G10\_T3\_L4\_C10\_dataset\_1.csv}\\\textcolor{DescText}{Data: 60 delivery-time observations, split evenly among North\_A, South\_B, and East\_C routes.}}, filenode, text width=16.5cm]
[{\code{G10\_T3\_L4\_C10\_dataset\_2.csv}\\\textcolor{DescText}{Data: 60 weekly screen-time observations, split among Grade 10 STEM, arts, and sport groups.}}, filenode, text width=16.5cm]
[{\code{G10\_T3\_L4\_C10\_dataset\_3.csv}\\\textcolor{DescText}{Data: 60 resale-value observations, split among Plan\_A, Plan\_B, and Plan\_C.}}, filenode, text width=16.5cm]
[{\code{G11\_T3\_C1C7\_midterm.tex}\\\textcolor{DescText}{Student-facing midterm on a Grade 11 punctuality audit, assessing C1–C7 across two cohort benchmarks.}}, filenode, text width=16.5cm]
[{\code{G11\_T3\_C1C7\_midterm\_solution.tex}\\\textcolor{DescText}{Teacher material: full worked midterm key with explanations and plotted normal rejection regions.}}, filenode, text width=16.5cm]
[{\code{G11\_T3\_C1C7\_midterm\_solution\_minimal.tex}\\\textcolor{DescText}{Teacher material: condensed key for the same midterm, retaining calculations/graphs with less exposition.}}, filenode, text width=16.5cm]
[{\code{G11\_T3\_L1\_C1C2C3\_exam\_solution.tex}\\\textcolor{DescText}{Teacher material: model C1–C3 contexts plus guidance for six student-designed school-to-Bogotá problems.}}, filenode, text width=16.5cm]
[{\code{G11\_T3\_L1\_C1C2C3\_exam\_student.tex}\\\textcolor{DescText}{Student-facing C1–C3 assessment template requiring six balanced mean/proportion hypothesis designs.}}, filenode, text width=16.5cm]
[{\code{G11\_T3\_L1\_C4C5\_exam\_student.tex}\\\textcolor{DescText}{Student-facing C4–C5 assessment requiring four data-supported tests selected from six personal/local contexts.}}, filenode, text width=16.5cm]
[{\code{G11\_T3\_L1\_C6C7\_exam\_solution.tex}\\\textcolor{DescText}{Teacher material: worked C6 decision/error cases and C7 minimum-sample-size cases in economic contexts.}}, filenode, text width=16.5cm]
[{\code{G11\_T3\_L2\_C4C5\_exam\_solutions.tex}\\\textcolor{DescText}{Teacher material: concise five-problem C4–C5 key covering foundational, sensitivity, directional, and synthesis tests.}}, filenode, text width=16.5cm]
[{\code{G11\_T3\_L4\_C10\_dataset\_1.csv}\\\textcolor{DescText}{Data: 30 study-hours observations, 15 each for groups A and B.}}, filenode, text width=16.5cm]
[{\code{G11\_T3\_L4\_C10\_dataset\_2.csv}\\\textcolor{DescText}{Data: 30 delivery-time observations, 15 each for FastShip and CityExpress.}}, filenode, text width=16.5cm]
[{\code{G11\_T3\_L4\_C10\_dataset\_3.csv}\\\textcolor{DescText}{Data: 30 exercise-duration observations, 15 each for weekdays and weekends.}}, filenode, text width=16.5cm]
[{\code{G11\_T3\_L4\_C8C9C10\_activity.tex}\\\textcolor{DescText}{Student-facing C8–C10 brief for an HTML/CSS/JavaScript two-sample-testing page and written report.}}, filenode, text width=16.5cm]
[{\code{G11\_T3\_L4\_C8C9C10\_dataset\_1.csv}\\\textcolor{DescText}{Data: 24 numeric scores for Control and Treatment groups (dataset G11\_D1).}}, filenode, text width=16.5cm]
[{\code{G11\_T3\_L4\_C8C9C10\_dataset\_2.csv}\\\textcolor{DescText}{Data: 30 binary success outcomes for Method\_A and Method\_B (dataset G11\_D2).}}, filenode, text width=16.5cm]
[{\code{G11\_T3\_L4\_C8C9C10\_dataset\_3.csv}\\\textcolor{DescText}{Data: 24 time measurements for Baseline and New\_Process groups (dataset G11\_D3).}}, filenode, text width=16.5cm]
[{\code{G11\_T3\_L4\_C8C9\_example\_mean\_data.csv}\\\textcolor{DescText}{Example data: 12 paired columns of numeric observations for groups A and B.}}, filenode, text width=16.5cm]
[{\code{G11\_T3\_L4\_C8C9\_example\_proportion\_data.csv}\\\textcolor{DescText}{Example data: 20 student-level group labels and binary success outcomes.}}, filenode, text width=16.5cm]
]
\end{forest}
\end{center}
\end{landscape}
\begin{landscape}
\subsection{\code{3-catch\_ups/}}
\begin{center}
\begin{forest}
for tree={grow=east, parent anchor=east, child anchor=west, anchor=west, edge={-Latex,draw=TreeLine}, l sep=9mm, s sep=1.2mm, font=\sffamily\fontsize{6.2}{7.2}\selectfont},
[{\textbf{3-catch\_ups/}\\\textcolor{DescText}{Remedial or supplementary assessment material.}}, rootnode, text width=4.5cm
[{\code{suplementary\_homework.tex}\\\textcolor{DescText}{Student-facing catch-up homework: a full C1–C7 hypothesis-design/decision task using two cohort benchmarks.}}, filenode, text width=16.5cm]
]
\end{forest}
\end{center}
\end{landscape}
\begin{landscape}
\subsection{\code{3-exams/}}
\begin{center}
\begin{forest}
for tree={grow=east, parent anchor=east, child anchor=west, anchor=west, edge={-Latex,draw=TreeLine}, l sep=9mm, s sep=1.2mm, font=\sffamily\fontsize{6.2}{7.2}\selectfont},
[{\textbf{3-exams/}\\\textcolor{DescText}{Expanded worked examination resources, separate from the learning-evidence folder.}}, rootnode, text width=4.5cm
[{\code{G11\_T3\_L2\_exam\_solutions.tex}\\\textcolor{DescText}{Teacher material: eight worked C4–C5 proportion/mean tests with graphs and contextual interpretation.}}, filenode, text width=16.5cm]
]
\end{forest}
\end{center}
\end{landscape}

\code{STRUCTURE.md} is this inventory and is intentionally omitted from the tree.

\section{Resource types and important relationships}

\begin{itemize}
\item \textbf{LaTeX documents and shared configuration:} most practice and solution
  documents load \code{../../../preamble/preamble\_exams.tex} (some use a flexible
  \code{\textbackslash{}input@path}); the C6–C7 solution instead loads \code{preamble.tex}. These shared
  preambles live outside Term 3 and supply document formatting and commands.
  Several standalone assessments declare their own packages rather than loading
  a shared preamble.
\item \textbf{Images and generated graphics:} standalone assessment headers reference
  \code{../../../preamble/logo.png}, also outside this tree. Files using
  \code{\textbackslash{}IfFileExists} render a boxed “Logo” fallback, while three early learning
  evidence files reference the logo directly. Statistical diagrams are drawn
  from LaTeX/TikZ/PGFPlots code in the sources; no separate Term 3 image or media
  assets are stored here.
\item \textbf{Practice progression:} the C1–C3 activities introduce parameters,
  hypotheses, and tail choice; the C4 and C5 ladders apply proportion and mean
  z-tests; C6 interprets decisions and errors; and C7 derives minimum sample
  sizes. \code{G11\_T3\_summary.tex} describes this sequence, although some filenames
  cited inside it use \code{L1} where the present files use \code{L2}/\code{L3} (see
  \textbf{Requires verification}).
\item \textbf{Assessment variants:} the C1–C3 \code{exam\_student} file is the blank student
  template and \code{exam\_solution} adds model solutions/instructor guidance. The
  midterm has one student paper and two keys: a fully explained key and a
  deliberately shorter “minimal” key. The two Lesson 2 solution files are not
  duplicates: \code{2-learning\_evidences/...exam\_solutions.tex} is a concise
  five-problem set, while \code{3-exams/...exam\_solutions.tex} expands coverage to
  eight worked C4–C5 problems.
\item \textbf{Web, stylesheets, and JavaScript:} no \code{.html}, \code{.css}, or \code{.js} files are
  present. \code{G11\_T3\_L4\_C8C9C10\_activity.tex} is the specification that requires
  students to submit a main HTML file, linked stylesheet(s), JavaScript file(s),
  and a deployed-page link. It also links to external Git/GitHub Pages tutorials
  and three SharePoint datasets; those links are not local file dependencies.
\item \textbf{Local datasets:} the activity brief does not name or programmatically load
  any local CSV. The \code{G11\_T3\_L4\_C8C9C10\_dataset\_*} files nevertheless match its
  requested two-sample numerical/proportion work: score and time data support
  mean comparisons, while binary success data support a proportion comparison.
  The two \code{example\_*} files provide smaller demonstration layouts. The three
  \code{G11\_T3\_L4\_C10\_dataset\_*} files are additional two-group numeric datasets.
\item \textbf{Downloadable documents and builds:} the directory contains editable \code{.tex}
  sources but no compiled PDFs or other downloadable binary documents. It also
  contains no build configuration, generated LaTeX auxiliaries, or dependency
  directories.
\end{itemize}

\section{Requires verification}
\begin{verificationbox}

\begin{itemize}
\item The three \code{G10\_T3\_L4\_C10\_dataset\_*} files contain explicitly Grade 10 group
  names/data but are stored in Grade 11 Term 3. Their contents are clear, but
  whether their placement is intentional cannot be established from local
  references.
\item \code{G11\_T3\_summary.tex} names C4–C7 practice files as Lesson 1 resources; the
  actual files are named Lesson 2 (C4/C5) and Lesson 3 (C6/C7). The summary's
  conceptual descriptions match those files, but the stale naming should be
  confirmed before treating it as a literal dependency map.
\end{itemize}

\end{verificationbox}
\section{Inventory scope}

This inventory documents \textbf{32 pre-existing files} and \textbf{4 directories} below
the Term 3 root. No transient or generated item was present to exclude. Per the
requested policy, this newly created inventory is not counted or listed in its
own tree.
\end{document}
